% Keep the chapter's reading order and drawing styles in one place.
\usetikzlibrary{patterns}
\colorlet{indexBlue}{ThesisBlue}
\colorlet{indexOrange}{ThesisOrange}
\colorlet{indexGreen}{ThesisGreen}
\tikzset{
  index cell/.style={draw=black!55,minimum height=.52cm,
    minimum width=.72cm,inner sep=3pt,font=\footnotesize,align=center},
  index blue/.style={index cell,draw=indexBlue,fill=indexBlue!12},
  index orange/.style={index cell,draw=indexOrange,fill=indexOrange!14},
  index green/.style={index cell,draw=indexGreen,fill=indexGreen!12},
  index arrow/.style={-{Latex[length=1.8mm]},thick}
}
\lstdefinestyle{bioseq-index-sql}{style=bioseq-search-sql}
\DeclareCodeRepository{minimap2index}{https://github.com/lh3/minimap2}[79c9cc186b95f50bd899f69b48eba995ced810c6]
\DeclareCodeRepository{mmseqsindex}{https://github.com/soedinglab/MMseqs2}[c77b5afa910bec52c784566e378cb6ebd3d0d453]
\raggedbottom
\section{Keep the lookup work between queries}
\label{sec:index-storage-role}

In previous chapter~\ref{sec:pg-storage-role} I describe how columnar and sequence storage layers are implemented in BioseqDB. There's however also aspect for storing indexes for sequences.
We would like to be able to store computed minimizers (seeds/sketches)~\ref{lst:minimap-minimizer-cpp-decl} in a way that enables performant computations of alignments.
They should be computed once when sequence is loaded from file or updated and then reused for subsequent select queries.

I want to discuss existing aligners' indexing first.
\section{What existing sequence indexes store}
\label{sec:index-minimap2-comparison}

\subsection{Indexing: minimap2}

After computing minimizers they need to be grouped by their hash, minimap2 choose to use the lowest $b$ hash bits (by default $b=14$) to choose one of $2^b$ buckets~\autocite{minimap2_paper,minimap2_software}.
Minimap2 builds a hash table per bucket whose keys are the remaining hash bits. This buckets with hasmaps are serialized into \path{.mmi} file.
I performed in-depth analysis of how the \path{.mmi} format is specified based on available Github old issues~\autocite{minimap2_issue_820,minimap2_mmi_gist} and original code from minimap2's \texttt{mm\_idx\_dump} and \texttt{mm\_idx\_load} functions.
\noindent
\coderef{minimap2index/index.c}{Writing and loading an index part}{mm_idx_dump, mm_idx_load}{475}

The index file is concatenation of blocks presented in the Figure~\ref{fig:index-mmi-layout}:
\begin{itemize}
  \item \textbf{Magic bytes} The bytes are \texttt{MMI} followed by
    \texttt{0x02} (\texttt{MM\_IDX\_MAGIC}).
  \item \textbf{Metadata} Minimap2 saves $w$, the $k$-mer
    length $k$, used bitcount $b$, the number of sequences
    $n_{\mathrm{seq}}$ and a flag word (for various minimap2-specific bitflags)
  \item \textbf{Sequence table} For each of the $n_{\mathrm{seq}}$ sequences
    the file stores 1-byte name length $\ell$ followed by $\ell$ ASCII characters representing the sequence name (header taken from original FASTA file), and
    a 32-bit sequence length
  \item \textbf{Buckets} The buckets store the hash-map, where key is remaining hash and bit indicating if there's single value for that key or not. 
    The hash-map is stored as two components however. 
    First is the value list $p$ which is used when there are more than 1 value per key.
    Second is the hash array itself, which stores key/value pairs.
    If there are more than one value per key (last bit of key set to 1), then the value represents offset and count into the value list $p$
    The values, either as values in the key/value pairs or entries in $p$ store sequence id, position and strand bit.
  \item \textbf{Bit-acked bases (optional)} If the \texttt{MM\_I\_NO\_SEQ} flag
    is unset manually by user, bit-packed reference bases are saved at the end.
\label{lst:high-level-mmi-index-format-description}
\end{itemize}

The depiction of the index file layout is presented in the Figure~\ref{fig:index-mmi-layout} below.

\begin{figure}[htbp]
\centering
\begin{tikzpicture}[
  x=1cm,y=1cm,font=\footnotesize,
  bitfill/.style={draw=black!60,line width=0.4pt},
  lab/.style={align=center,font=\scriptsize,inner sep=0pt},
  zoom/.style={line width=0.9pt, line join=miter, line cap=rect,
    -{Latex[length=2.4mm,width=1.7mm]}}
]
  % (b) indent 1.05, (c)=(d) indent 2.10.
  % (b) ends at 8.55 so the green rail can pass at x=9.35.

  % ----- (a) y: boxes 14.20--14.95 -----
  \node[anchor=south west,font=\bfseries\footnotesize] at (0,15.25)
    {(a) Byte order of one \texttt{.mmi} part};
  \fill[indexBlue!14] (0,14.20) rectangle (1.70,14.95);
  \draw[black!55] (0,14.20) rectangle (1.70,14.95);
  \node[lab] at (0.85,14.575) {magic\\4\,B};
  \fill[indexBlue!14] (1.70,14.20) rectangle (5.20,14.95);
  \draw[black!55] (1.70,14.20) rectangle (5.20,14.95);
  \node[lab] at (3.45,14.575) {$w$, $k$, $b$, $n_{\mathrm{seq}}$, flag\\$5\times 4$\,B};
  \fill[indexOrange!16] (5.20,14.20) rectangle (7.80,14.95);
  \draw[black!55] (5.20,14.20) rectangle (7.80,14.95);
  \node[lab] at (6.50,14.575) {sequence table\\$n_{\mathrm{seq}}$ records};
  \fill[indexGreen!16] (7.80,14.20) rectangle (10.90,14.95);
  \draw[black!55] (7.80,14.20) rectangle (10.90,14.95);
  \node[lab] at (9.35,14.575) {buckets\\$2^{b}$ in order};
  \fill[black!8] (10.90,14.20) rectangle (12.80,14.95);
  \draw[black!55] (10.90,14.20) rectangle (12.80,14.95);
  \node[lab] at (11.85,14.575) {packed $S$\\if stored};

  % ----- (b) boxes 11.45--12.20, title below orange rail -----
  \node[anchor=south west,font=\bfseries\footnotesize] at (1.05,12.50)
    {(b) One sequence record};
  \fill[indexOrange!16] (1.05,11.45) rectangle (2.95,12.20);
  \draw[black!55] (1.05,11.45) rectangle (2.95,12.20);
  \node[lab] at (2.00,11.825) {name length $\ell$\\1\,B};
  \fill[indexOrange!16] (2.95,11.45) rectangle (6.45,12.20);
  \draw[black!55] (2.95,11.45) rectangle (6.45,12.20);
  \node[lab] at (4.70,11.825) {name (ASCII)\\$\ell$\,B};
  \fill[indexOrange!16] (6.45,11.45) rectangle (8.55,12.20);
  \draw[black!55] (6.45,11.45) rectangle (8.55,12.20);
  \node[lab] at (7.50,11.825) {seq.\ length\\4\,B};

  % ----- (c) boxes 8.75--9.50, title below green rail -----
  \node[anchor=south west,font=\bfseries\footnotesize] at (2.10,9.80)
    {(c) One bucket};
  \fill[indexGreen!16] (2.10,8.75) rectangle (3.70,9.50);
  \draw[black!55] (2.10,8.75) rectangle (3.70,9.50);
  \node[lab] at (2.90,9.125) {$n$\\4\,B};
  \fill[indexOrange!16] (3.70,8.75) rectangle (7.10,9.50);
  \draw[black!55] (3.70,8.75) rectangle (7.10,9.50);
  \node[lab] at (5.40,9.125) {value list $p$\\$n\times 8$\,B};
  \fill[indexGreen!16] (7.10,8.75) rectangle (8.60,9.50);
  \draw[black!55] (7.10,8.75) rectangle (8.60,9.50);
  \node[lab] at (7.85,9.125) {$s$\\4\,B};
  \fill[indexBlue!14] (8.60,8.75) rectangle (12.80,9.50);
  \draw[black!55] (8.60,8.75) rectangle (12.80,9.50);
  \node[lab] at (10.70,9.125) {$s$ key/value pairs\\$s\times(8{+}8)$\,B};

  % ----- (d) -----
  \node[anchor=south west,font=\bfseries\footnotesize] at (2.10,7.20)
    {(d) 64-bit words, high bits on the left ($b{=}14$)};

  \node[anchor=south west,font=\scriptsize] at (2.10,6.80) {Hash-table key};
  \fill[black!10] (2.10,6.05) rectangle (5.6109,6.65);
  \draw[bitfill] (2.10,6.05) rectangle (5.6109,6.65);
  \node[lab] at (3.86,6.35) {unused, 21 bits};
  \fill[indexBlue!18] (5.6109,6.05) rectangle (12.6328,6.65);
  \draw[indexBlue] (5.6109,6.05) rectangle (12.6328,6.65);
  \node[lab] at (9.12,6.35) {remaining hash $(\mathit{hash}\gg b)$, 42 bits};
  \fill[indexGreen!35] (12.6328,6.05) rectangle (12.80,6.65);
  \draw[indexGreen,very thick] (12.6328,6.05) rectangle (12.80,6.65);

  \node[anchor=south west,font=\scriptsize] at (2.10,5.50)
    {Hash-map value for single value / $p$ element};
  \fill[indexBlue!18] (2.10,4.70) rectangle (7.45,5.30);
  \draw[indexBlue] (2.10,4.70) rectangle (7.45,5.30);
  \node[lab] at (4.78,5.00) {sequence id, 32 bits};
  \fill[indexOrange!20] (7.45,4.70) rectangle (12.6328,5.30);
  \draw[indexOrange] (7.45,4.70) rectangle (12.6328,5.30);
  \node[lab] at (10.04,5.00) {last-base position, 31 bits};
  \fill[indexGreen!35] (12.6328,4.70) rectangle (12.80,5.30);
  \draw[indexGreen,very thick] (12.6328,4.70) rectangle (12.80,5.30);

  \node[anchor=south west,font=\scriptsize] at (2.10,4.20)
    {Hash-map value for non-single value};
  \fill[indexOrange!20] (2.10,3.40) rectangle (7.45,4.00);
  \draw[indexOrange] (2.10,3.40) rectangle (7.45,4.00);
  \node[lab] at (4.78,3.70) {offset into $p$, 32 bits};
  \fill[indexBlue!18] (7.45,3.40) rectangle (12.80,4.00);
  \draw[indexBlue] (7.45,3.40) rectangle (12.80,4.00);
  \node[lab] at (10.12,3.70) {count, 32 bits};

  \fill[indexGreen!35] (2.10,2.75) rectangle (2.38,3.03);
  \draw[indexGreen] (2.10,2.75) rectangle (2.38,3.03);
  \node[anchor=west,font=\scriptsize,text=black!55] at (2.52,2.89)
    {1-bit field: either means strand (for hash-map value) or means that hash-map value is single (for hash-map key)};

  % Orange (a seq -> b): H at mid(14.20, 12.20)=13.20. Spine 0.22.
  \draw[zoom,indexOrange]
    (6.50,14.20) -- (6.50,13.20) -- (0.22,13.20) -- (0.22,11.825) -- (1.05,11.825);
  % Green (a buckets -> c): H at mid(11.45, 9.50)=10.475. Spine 0.55.
  \draw[zoom,indexGreen]
    (9.35,14.20) -- (9.35,10.475) -- (0.55,10.475) -- (0.55,9.125) -- (2.10,9.125);
  % Blue (pairs -> key): H 7.95, below the p rail. Spine 1.25.
  \draw[zoom,indexBlue]
    (9.80,8.75) -- (9.80,7.95) -- (1.25,7.95) -- (1.25,6.35) -- (2.10,6.35);
  % Orange (p -> occurrence): H 8.40, above the blue rails. Spine 0.40.
  \draw[zoom,indexOrange]
    (5.40,8.75) -- (5.40,8.40) -- (0.40,8.40) -- (0.40,5.00) -- (2.10,5.00);
  % Blue (pairs -> value): right rail 13.60, enter from the right.
  \draw[zoom,indexBlue]
    (11.80,8.75) -- (11.80,7.95) -- (13.60,7.95) -- (13.60,3.70) -- (12.80,3.70);
  \draw[zoom,indexBlue]
    (11.80,8.75) -- (11.80,7.95) -- (13.60,7.95) -- (13.60,5.00) -- (12.80,5.00);
\end{tikzpicture}
\caption[Interpretation of one \texttt{.mmi} index]{%
Interpretation of one \texttt{.mmi} index.
For very large datasets index files are split into multiple parts where each of them follows the same format.
Minimap2 needs to store sequence ID (RID~\ref{lst:minimap-minimizer-cpp-decl}), because the index parts are generated for a set of sequences, not just one, so
each part needs to identify which reference sequence the minimizer comes from.
Arrows represent how specific elements of arrays look like in close inspection.
Panel~(a) depicts high-level structure of the index~\ref{lst:high-level-mmi-index-format-description}.
Panel~(b) presents records storing sequence information.
Panel~(c) shows each bucket of minimizers, which stores its own hash-map via two lists.
Panel~(d) shows how elements of $p$ looks like and key/values in the hash-map list.
Hash map key is always remaining 42 bits with bit indicating that there's single value for that key.
If the bit is set, then value in the key/value pair is (sequence id, position, strand bit), if not, then value points via offset/count pair into the list $p$.
The list $p$ stores the same (sequence id, position, strand bit) values.}
\label{fig:index-mmi-layout}
\end{figure}

The idea of storing hash-map as two lists inside bucket is also presented in the Figure~\ref{fig:index-mmi-construction} below to allow easier understanding of the 
index concept.

\begin{figure}[htbp]
\centering
\begin{tikzpicture}[x=1cm,y=1cm,font=\footnotesize]
  \node[anchor=west,font=\bfseries] at (0,5.8) {Minimizer occurrences before grouping};
  \foreach \x/\h/\p in {1/13/17,3/21/24,5/10/29,7/21/30}
    \node[index blue,minimum width=1.6cm] at (\x,4.95) {$h=\h$\\$p=\p$};
  \node[anchor=west,align=left] at (9,4.95) {$b=2$: bucket $h\bmod4$\\remaining key: $\lfloor h/4\rfloor$};
  \foreach \x/\b in {1/0,4/1,7/2,10/3}
    \node[index cell,fill=black!3,minimum width=2.2cm] (b\b) at (\x,3.3) {bucket \b};
  \draw[index arrow,indexBlue] (1,4.42)--(4,3.6);
  \draw[index arrow,indexBlue] (3,4.42)--(4,3.6);
  \draw[index arrow,indexBlue] (7,4.42)--(4,3.6);
  \draw[index arrow,black!45] (5,4.42)--(7,3.6);
  \node[anchor=west,font=\bfseries] at (0,2.35) {Bucket 1 after sorting};
  \node[index blue,minimum width=3.6cm] (single) at (2.35,1.5)
    {remaining bits 3: $p=17$};
  \node[index orange,minimum width=3.6cm] (repeated) at (2.35,.5)
    {remaining bits 5: offset 0, count 2};
  \node[anchor=west] at (6.0,1.7) {separate position array};
  \node[index orange,minimum width=1.15cm] at (6.65,.5) {$p=24$};
  \node[index orange,minimum width=1.15cm] at (7.85,.5) {$p=30$};
  \node[index cell,dashed,minimum width=1.15cm] at (9.05,.5) {$\cdots$};
  \node at (6.65,1.08) {0}; \node at (7.85,1.08) {1};
  \draw[index arrow,indexOrange] (repeated.east)--(6.05,.5);
  \node[anchor=west,text=black!65] at (0,-.45)
    {A singleton needs no separate list. Repeated keys share a contiguous array.};
\end{tikzpicture}
\caption{A simplified minimap2 hash bucket: hashes 13 and 21 have the same
lowest two bits, but different remaining keys, so they land in bucket 1 together. The two occurrences of hash
21 are placed next to each other in the bucket's position array.}
\label{fig:index-mmi-construction}
\end{figure}

\subsection{Indexing: mmseqs2}

There are alternative approaches in literature to building an index.
MMseqs2 stores one index file that contains counts and offsets, and then writes larger index files 
that contain all occurence lists instead of bucketing them like minimap2 does.
It's worth noting that 2017 MMseqs2\autocite{Steinegger079681,mmseqs2_2017} paper uses spaced k-mers which wasn't the focus of this thesis, so I didn't 
investigate more how the alignment is precisely done.

MMseqs2 offers great performance in clsutering, which is different problem to alignment, however connected as it involves 
estimation of edit distances for large datasets. The software also allows sets of taxonomy operations and alignment-based searches as well.\autocite{mmseqs2_linclust_2018}\autocite{mmseqs2_taxonomy_2021}\autocite{mmseqs2_desktop_2019}\autocite{mmseqs2_gpu_2025}.
\noindent
\coderef{mmseqsindex/src/prefiltering/IndexBuilder.cpp}{Count and fill the MMseqs2 index}{IndexBuilder::fillDatabase}{56}

\noindent
\coderef{mmseqsindex/src/prefiltering/IndexTable.h}{Read a word's entry range}{getDBSeqList}{178}

\noindent
\coderef{mmseqsindex/src/prefiltering/IndexTable.h}{Keep distinct word codes within one target}{addSequence}{350}

I illustrate the idea of MMseqs2 index on a small example in the Figure~\ref{fig:index-mmseqs-inverted} below, where 
I just assume that we use 2-letter k-mers of natural language for simplicity.

\begin{figure}[htbp]
\centering
\begin{tikzpicture}[x=1cm,y=1cm,font=\footnotesize]
  \node[anchor=west,font=\bfseries] at (0,6.55) {(a) Read target words};
  \foreach \x/\id/\s in {2/1/ACD,6/2/CAD,10/3/ACE}
    \node[index blue,minimum width=2.8cm] at (\x,5.75) {target \id: \texttt{\s}};
  \node[anchor=west,font=\bfseries] at (0,4.8) {(b) Count entries, then sum counts to find list starts};
  \node[anchor=east] at (1.25,3.95) {word};
  \node[anchor=east] at (1.25,3.2) {count};
  \node[anchor=east] at (1.25,2.45) {offset};
  \foreach \x/\word/\n/\o in {2.25/AC/2/0,4.1/AD/1/2,5.95/CA/1/3,7.8/CD/1/4,9.65/CE/1/5} {
    \node[index blue,minimum width=1.55cm] at (\x,3.95) {\texttt{\word}};
    \node[index cell,minimum width=1.55cm] at (\x,3.2) {$\n$};
    \node[index green,minimum width=1.55cm] at (\x,2.45) {$\o$};
  }
  \node[index green,minimum width=1.25cm] at (11.25,2.45) {6};
  \node[above] at (11.25,2.78) {end};
  \node[anchor=west,font=\bfseries] at (0,1.5) {(c) Fill the reserved list ranges};
  \foreach \x/\id/\p/\i in {2.1/1/0/0,3.7/3/0/1,5.3/2/1/2,6.9/2/0/3,8.5/1/1/4,10.1/3/1/5} {
    \node[index cell,minimum width=1.48cm] at (\x,.5) {$(\id,\p)$};
    \node[above] at (\x,.81) {$\i$};
  }
  \draw[indexBlue,very thick] (1.3,.12)--(4.5,.12);
  \node[below,text=indexBlue] at (2.9,.12) {AC: entries $[0,2)$};
  \node[anchor=west,align=left] at (6,-.15)
    {Entry = (target ID, word start).\\Positions start at zero.};
\end{tikzpicture}
\caption{A small example of MMseqs2's list construction, where word/count/offsets triplets are stored in a single file - practically only indices are stored there, while the actual (very large) lists of entries are stored separately.}
\label{fig:index-mmseqs-inverted}
\end{figure}

\section{BioseqDB indexing abstractions}
\label{sec:index-abstraction-layers}

Indexes are stored separately in BioseqDB~\ref{lst:types-of-storage-entities-in-bioseqdb} to all internal sequence fragmented tables~\ref{fig:storage-length-buckets} 
and they are separate entities to the usual indexes in Postgres which provide generic access paths to the tables.

There is thin layer of abstraction over the indexes, so in the future implementation of other types of indexes is easy. The storage abstraction 
requires implementer only to specify the type and depending on it we only need to provide logic to write the index chunks.
One problem with current state of BioseqDB is that the index reading is not abstracted, so we encapsulate details of implementation, but they leak into pipeline (as the pipeline 
currently needs to understand how we store minimizers).

\info{Seed abstraction}{In the future I want to improve curent implementation to use abstract notion of "seed". I already extensively use templates, so 
introducing new concept to represents types of seeds isn't a problem. This abstraction will allow to fully decouple indexing from the pipeline.
Pipeline stages then could automatically detect how to compute certain seeds (like minimizers) and how to read and store them.
Current BioseqDB implementation only uses minimizers, so I decided not to decicate my time to this idea.
}

I assume here that index store \emph{chunks} which are separate abstraction to the chunks in columnar tables.
One chunk consists of bytes the index implementation needs to write or store and I assume the chunk cannot be larger than $960MB$ (due to Postgre's TOASTing limits). 
The chunks are stored in internal tables that does not use columnar storage, they're ordinary heap tables - I decided that columnar storage in this case won't give any advantage (although I didn't A/B test this approach).
Additionally, index tables are create \texttt{UNLOGGED} to avoid unnecessary logging and to reduce the amount of data written to the disk.

There are several supported methods of storing index:

\begin{itemize}
  \item \textbf{\texttt{CHUNKED\_PER\_FRAGMENT}}.
    One fragment~\ref{info:sequence-fragments} of a sequence can have multiple chunks in the index.
    % They are grouped by a \emph{sub-key}, an extra number chosen by the index.
    % For \texttt{minimizers\_per\_fragment} that number is the fragment, so one fragment can fill several chunks.
    % \texttt{minimizers\_per\_sequence} uses this same layout.
    % Its sub-key is the number of one hash group, not a fragment number.
  \item \textbf{\texttt{INLINE\_PER\_SEQUENCE}}.
    One sequence has a single chunk that is stored as Postgres varlena. This means that entire data stored for one sequence cannot exceed $960MB$ if that happens BioseqDB will error.
    \texttt{symbol\_counts} is BiosqeDB index that uses this layout. It stores how many times each symbol occurs in that sequence and the payload is so small that it nicely fits into one small row.
  \item \textbf{\texttt{CHUNKED\_PER\_TABLE}}.
    This is special layout that in its implementation exposes only simple methods to access random chunks of the table. The idea is that some indexing structures may want 
    to implement custom strategies, for example trees of sequences or graphs of fragments that do not fit into other representations.
    I don't use this layout at the moment, but it's supported to provide extension point in the future. 
\end{itemize}

Figure~\ref{fig:index-storage-layouts} shows what three of these layouts are representing.

\begin{figure}[htbp]
\centering
\begin{tikzpicture}[
  x=1cm,y=1cm,font=\footnotesize,
  lab/.style={align=center,font=\scriptsize,inner sep=1pt},
  arr/.style={line width=0.85pt,line join=miter,line cap=rect,
    -{Latex[length=2.4mm,width=1.7mm]}},
  stem/.style={line width=0.85pt,line join=miter,line cap=rect}
]
  \node[anchor=south west,font=\bfseries\footnotesize] at (0.85,12.22)
    {(a) User table};
  \draw[indexGreen,line width=1pt,fill=black!3]
    (0.85,9.35) rectangle (6.35,12.05);
  \node[lab,anchor=west] at (1.05,11.68) {id};
  \node[lab,anchor=west] at (2.45,11.68) {sequence locator};
  \draw[black!25] (0.85,11.42) -- (6.35,11.42);
  \fill[black!5] (1.05,10.72) rectangle (2.05,11.22);
  \draw[black!40] (1.05,10.72) rectangle (2.05,11.22);
  \node[lab] at (1.55,10.97) {2};
  \fill[black!5] (2.25,10.72) rectangle (6.15,11.22);
  \draw[black!40] (2.25,10.72) rectangle (6.15,11.22);
  \node[lab] at (4.20,10.97) {locator of sequence 2};
  \fill[black!5] (1.05,9.58) rectangle (2.05,10.28);
  \draw[black!40] (1.05,9.58) rectangle (2.05,10.28);
  \node[lab] at (1.55,9.93) {1};
  \fill[indexOrange!22] (2.25,9.58) rectangle (6.15,10.28);
  \draw[indexOrange,line width=1pt] (2.25,9.58) rectangle (6.15,10.28);
  \node[lab] at (4.20,9.93) {locator of sequence 1};

  % ---------- (b) fragments ----------
  \node[anchor=south west,font=\bfseries\footnotesize] at (7.15,12.22)
    {(b) Fragment rows of sequence 1};
  \draw[black!55,fill=black!3] (7.15,9.35) rectangle (15.05,12.05);
  \fill[black!5] (7.40,9.72) rectangle (9.55,11.55);
  \draw[black!40] (7.40,9.72) rectangle (9.55,11.55);
  \node[lab] at (8.475,10.64) {fragment 0};
  \fill[indexBlue!18] (9.80,9.72) rectangle (12.35,11.55);
  \draw[indexBlue,line width=1pt] (9.80,9.72) rectangle (12.35,11.55);
  \node[lab] at (11.075,10.64) {fragment 1};
  \fill[black!5] (12.60,9.72) rectangle (14.80,11.55);
  \draw[black!40] (12.60,9.72) rectangle (14.80,11.55);
  \node[lab] at (13.70,10.64) {fragment 2};

  % ---------- shared row geometry ----------
  % bar i top = 4.28 - i*0.56, height 0.46, center = 4.05 - i*0.56

  % ---------- (c) table chunks: all 5 rows ----------
  \draw[indexGreen!70,fill=indexGreen!5] (1.05,1.28) rectangle (4.85,5.62);
  \node[lab,font=\bfseries\scriptsize] at (2.95,5.22) {\texttt{CHUNKED\_PER\_TABLE}};
  \node[lab] at (2.95,4.72) {chunks of the whole table};
  \foreach \i/\name in {0/{chunk 0},1/{chunk 1},2/{chunk 2},3/{chunk 3},4/{chunk 4}} {
    \fill[indexGreen!20] (1.22,{4.28-\i*0.56}) rectangle (4.68,{3.82-\i*0.56});
    \draw[indexGreen] (1.22,{4.28-\i*0.56}) rectangle (4.68,{3.82-\i*0.56});
    \node[lab] at (2.95,{4.05-\i*0.56}) {\name};
  }

  % ---------- (d) inline: 5 rows, one belongs to sequence 1 ----------
  \draw[black!45,fill=indexOrange!4] (6.05,1.28) rectangle (9.85,5.62);
  \node[lab,font=\bfseries\scriptsize] at (7.95,5.22) {\texttt{INLINE\_PER\_SEQUENCE}};
  \node[lab] at (7.95,4.72) {one row for each sequence};
  \foreach \i/\name/\on in {0/{sequence 2}/0,1/{sequence 1}/1,2/{other sequence}/0,3/{other sequence}/0,4/{other sequence}/0} {
    \ifnum\on=1
      \fill[indexOrange!24] (6.22,{4.28-\i*0.56}) rectangle (9.68,{3.82-\i*0.56});
      \draw[indexOrange,line width=0.9pt] (6.22,{4.28-\i*0.56}) rectangle (9.68,{3.82-\i*0.56});
    \else
      \fill[black!4] (6.22,{4.28-\i*0.56}) rectangle (9.68,{3.82-\i*0.56});
      \draw[black!35] (6.22,{4.28-\i*0.56}) rectangle (9.68,{3.82-\i*0.56});
    \fi
    \node[lab] at (7.95,{4.05-\i*0.56}) {\name};
  }

  % ---------- (e) fragment chunks: 3 of 5 rows are fragment 1 ----------
  \draw[black!45,fill=indexBlue!4] (10.85,1.28) rectangle (15.05,5.62);
  \node[lab,font=\bfseries\scriptsize] at (12.95,5.22) {\texttt{CHUNKED\_PER\_FRAGMENT}};
  \node[lab] at (12.95,4.72) {chunks of fragment 1};
  \foreach \i/\name/\on in {
      0/{fragment 0}/0,
      1/{fragment 1, chunk 0}/1,
      2/{fragment 1, chunk 1}/1,
      3/{fragment 1, chunk 2}/1,
      4/{fragment 2}/0} {
    \ifnum\on=1
      \fill[indexBlue!18] (11.02,{4.28-\i*0.56}) rectangle (14.88,{3.82-\i*0.56});
      \draw[indexBlue] (11.02,{4.28-\i*0.56}) rectangle (14.88,{3.82-\i*0.56});
    \else
      \fill[black!4] (11.02,{4.28-\i*0.56}) rectangle (14.88,{3.82-\i*0.56});
      \draw[black!35] (11.02,{4.28-\i*0.56}) rectangle (14.88,{3.82-\i*0.56});
    \fi
    \node[lab] at (12.95,{4.05-\i*0.56}) {\name};
  }

  % Green: open brace on the whole user table, then one branch into every chunk row.
  % End caps stop short of the table border so they do not form an extra column.
  \draw[stem,indexGreen] (0.62,11.92) -- (0.42,11.92) -- (0.42,10.70)
    -- (0.28,10.70) -- (0.28,1.81);
  \draw[stem,indexGreen] (0.62,9.48) -- (0.42,9.48) -- (0.42,10.70);
  \foreach \i in {0,1,2,3,4} {
    \draw[arr,indexGreen] (0.28,{4.05-\i*0.56}) -- (1.22,{4.05-\i*0.56});
  }

  % Orange: locator cell -> the one sequence-1 row (bar 1, center 3.49).
  \draw[stem,indexOrange] (4.20,9.58) -- (4.20,8.55) -- (5.45,8.55) -- (5.45,3.49);
  \draw[arr,indexOrange] (5.45,3.49) -- (6.22,3.49);

  % Blue: leave fragment 1 itself, then split into its three chunk rows.
  \draw[stem,indexBlue] (11.075,9.72) -- (11.075,7.55) -- (10.22,7.55) -- (10.22,2.37);
  \foreach \i in {1,2,3} {
    \draw[arr,indexBlue] (10.22,{4.05-\i*0.56}) -- (11.02,{4.05-\i*0.56});
  }
\end{tikzpicture}
\caption[Logical concepts of the three index storage layouts]{Each index is drawn as five rows.
The green brace represents the fact that the index structure represents all sequences stored in the table.
The orange arrow shows that one locator corresponds to exactly one chunk (one row in the index table).
The blue arrow shows that one fragment can have multiple chunks.
In case of \textbf{\texttt{CHUNKED\_PER\_TABLE}} we only store chunk number and for \textbf{\texttt{CHUNKED\_PER\_FRAGMENT}} I store fragment's number and chunk number.
}
\label{fig:index-storage-layouts}
\end{figure}

% Keep the schema introduction, caption and SQL together across page breaks.
\noindent\begin{minipage}{\linewidth}
In reality the fragment's number can be viewed as a generic sub-key. This distorts the original meaning of \textbf{\texttt{CHUNKED\_PER\_FRAGMENT}}, because in relity it's just 
"per-sequence" but with extra key: $\langle$query\_key,\\sub\_key,\\chunk\_no$\rangle$, while \textbf{\texttt{CHUNKED\_PER\_TABLE}} has $\langle$chunk\_no$\rangle$.
This idea is depicted below by real SQL schema of internal storage:

\begin{lstlisting}[style=bioseq-index-sql,
  caption={One heap row of a per-fragment index.},
  label={lst:chunked-per-fragment-row}]
CREATE TABLE bioseq_internal.seqidx_mpf_... (
  query_key   BIGINT NOT NULL, -- which sequence in this shard
  sub_key     BIGINT NOT NULL, -- which item of that sequence (can be fragment number but doesn't practically have to)
  -- Then we store byte array defined by index implementation but split across chunks:
  chunk_no    INT    NOT NULL, -- 0, 1, 2, ... of that item's byte string
  chunk_bytes INT    NOT NULL, -- how many bytes this chunk holds
  blob        BYTEA  NOT NULL, -- the bytes of this chunk
  PRIMARY KEY (query_key, sub_key, chunk_no)
);
\end{lstlisting}
\end{minipage}

The payload of an index is one byte array.
\texttt{CHUNKED\_PER\_FRAGMENT} stores that array as rows of Listing~\ref{lst:chunked-per-fragment-row}.
When the array is longer than one chunk, it is divided into chunk 0, chunk 1, and the following chunks.
Those rows keep the same \texttt{query\_key} and the same \texttt{sub\_key}.
A reader orders them by \texttt{chunk\_no} and concatenates \texttt{blob}.

\info{Storage details}{
\label{info:index-storage-details}
  It's important to note that for minimizers per fragment MPF the sub-key number is the fragment number.
I use the same column to store another number when the index is not split by fragment.
Technically this means that \textbf{\texttt{CHUNKED\_PER\_FRAGMENT}} should have its name changed to something more appropriate.
}

In code \texttt{pg\_supported\_indexes} is a special templated that connects those implementation classes to the registered indexes. This is frequent practice used by me - encoding 
implementation vs. abstraction pairs as tuples in type lists created using template metaprogramming, so if extension is required we only need to write new type and add it into the 
\texttt{pg\_supported\_indexes} template list. Checks if implementation is valid are done via C++26 concepts~\ref{lst:mpf-index-storage-check}:

\begin{BioseqCppSnippet}[caption={Minimizers per fragment require fragment-chunk storage.},
  label={lst:mpf-index-storage-check}]
template <class p_storage>
struct minimizers_per_fragment_index {
  static_assert(is_index_chunked_storage_for<
      p_storage, sequence_index_layout::fragment_chunks>);
  static constexpr sequence_index_layout index_type =
      sequence_index_layout::fragment_chunks;
  using storage_type = p_storage;
};
\end{BioseqCppSnippet}

\section{BioseqDB supported indexing}
\label{sec:index-logical-layouts}

There are 3 implementations of indexes inside BioseqDB:
\begin{itemize}
  \item \textbf{\texttt{minimizers\_per\_fragment}} (shortened to \texttt{mpf}) stores minimizers calculated per one fragment and uses \textbf{\texttt{CHUNKED\_PER\_FRAGMENT}} layout.
  \item \textbf{\texttt{minimizers\_per\_sequence}} (shortened to \texttt{mps}) stores minimizers calculated per sequence, but as the number of minimizers grows
    they would not fit in a single chunk, thus one may consider using \textbf{\texttt{CHUNKED\_PER\_TABLE}}.
    However, as mentioned, the sub-key~\ref{info:index-storage-details} is a number, so in MPS I use this number as a hash-group id, not as a hash.
    The rows stay in the \textbf{\texttt{CHUNKED\_PER\_FRAGMENT}} layout even though
    these numbers are not fragment numbers, as seen on the figure below~\ref{fig:index-mpf-mps}.
  \item \textbf{\texttt{symbol\_counts}}: stores counter of occurences of each symbol in a sequence, uses \textbf{\texttt{INLINE\_PER\_SEQUENCE}} layout, because the counter struct is very small,
    so Postgres inlines the varlena into the row.
\label{lst:types-of-indexes-in-bioseqdb}
\end{itemize}




The main idea behind MPF and MPS was:
\begin{itemize}
  \item \textbf{MPF}: I can keep small hashmaps separated so random queries during seeding stage, where we need to find common minimizers between query and target, can hit
  CPU caches more often as the accesses are localized.
  \item \textbf{MPS}: The most frequent minimizers are trimmed by minimap2 and MPF trimming is done per fragment, so it's less reliable and can remove minimizers 
  that are locally frequent, but globally rare reducing total number of minimizers and speeding up processing at the cost of potentially missing some matches. Additionally,
  MPF needs to perform multiple hash lookups if we want to determine all locations of minimizer's occurences so there's also performance penalty. This penalty can be 
  also mitigated by the fact that MPF is easier to parallelize.
  Thus we can see there are several important trade-offs and it's not clear if the implementation of MPS should be more or less performant in theory.
  In case of BioseqDB as it would be very inneficient to store all minimizers individually, I decided to still generate minimizers per fragment, but store them and merge later
  during process that can potentially spill to disk (if the entire index fits into RAM then no spillage occurs, but if it doesn't there are heavy performance penalties).
  Additionally I must note that MPS uses the trick described here~\ref{info:index-storage-details}, so the sub-key is a hash-group number, not fragment number.
  This implies that if the query is short there will be not many groups to check, and as each group is held as heap rows, then MPS can be very fast, because it reads only the list of one group, not the entire index.
\end{itemize}

The figure~\ref{fig:index-mpf-mps} below shows 6 minimizer occurrences in a sequence, and panel~(c) shows the byte strings that store them.
Panels~(a) and~(b) show the same six minimizers in two groupings.
Panel~(c)and later illustrate the encoding of the strctures in binary strings (every rectangle onward is the binary string \texttt{blob} of one heap row from Listing~\ref{lst:chunked-per-fragment-row}), where 
the label on the left is that row's \texttt{query\_key} and \texttt{sub\_key}.
If a string is longer than one chunk, later chunks continue it under the same two numbers.

\subsubsection{Binary coding of non-grouped MPF}

There are two types of MPF index depedning on BioseqDB setting flag: \path{bioseq.aligner.grouped_soa_generation.enabled}.
The non-grouped version is depicted in panel~(c). The sub-key is just the fragment number. The entry intself contains two arrays similar to what minimap2~\ref{fig:index-mmi-layout} does.
The first array with \path{controls} contains hash value, start position in \path{positions} array and count of how many positiosn have that same hash.
\path{positions} just represent minimizer positions.
Usually in headers there's additional metadata and some of it is redundant as the exact shapes evolved during my implementation. 
There are versions and special magic values to potentially error if index was corrupted.

\subsubsection{Binary coding of grouped MPF}

I tried to improve MPF by doing something more sophisticated. The MPF can changed format based on grouped generation setting flag.
Panel~(d) shows how grouped MPF looks like:
\begin{itemize}
    \item This is similar to MPF buit instead of storing array of structures (AoS) we store structure of arrays (SoA) so hash, start, and count are now three arrays in a row.
    \item I use special bit i.e $2^{61}$ to store positions array separately in separate row.
    \item There also is one special directory row using unique \path{sub\_key} value. 
    The directory is list of $128B$ entries mapping fragment number to specific \path{sub\_key} value.
    \item Each directory row $128B$ entry constains useful metadata like min(hash), max(hash) so some rows doesn't have to be read at all if they don't contain 
    hash values we are interested in.
    \item Using this \path{sub\_key}$\leftrightarrow$\path{fragment} mapping allows to potentially split fragment hash/start/count entry into multiple ones. 
    However this is not used in the code currently and effectively \path{sub_key} $=$ \path{fragment} for all entries.
    \item The directory item also contains information that is redundant for example \path{controls\_first} is number of first chunk storing hash/start/count binary string.
    This is always $0$ as the binary strings are split into chunks, which are zero-indexed. The same logic applies to other fields and the entry size 
    could be massively reduced.
    \item Grouped-MPF is very generic format and it also supports variable-width entries. The width of hash/start/count/positions can be adjusted independently for example 
    for very short sequences we could use $2B$ \path{position} array and for very low minimizer density we could use $2B$ for \path{start\_array} or 
    for very unrepetitive sequences where unique minimizers are very common, we could use $2B$ for \path{count\_array}.
    Practically, current code uses currently same widths for most of the fields and hash is adjusted depending on $k$ and
    sequence length from minimizer type~\ref{fig:minimizer-representation-layout}.
\end{itemize}

On the figure Panel~(d) I didn't depict all fields in the MPF grouped directory as not all fields are relevant.
Table~\ref{tab:grouped-mpf-directory-entry} below lists all fields that are stored:
\begin{table}[htbp]
\centering
\small
\caption{Fields of one 128-byte grouped MPF directory entry.}
\label{tab:grouped-mpf-directory-entry}
\begin{tabularx}{\linewidth}{@{}llX@{}}
\toprule
\textbf{Bytes} & \textbf{Size} & \textbf{Field}\\
\midrule
0--3 & 4 & partition number (this is an index of entry in the directory row's array)\\
4--7 & 4 & fragment\\
8--15 & 8 & min(hash)\\
16--23 & 8 & max(hash)\\
24--31 & 8 & sub-key (in theory was meant to use to allow dividing entry for fragment into multiple rows, but currently it's always equal to frament number)\\
32--35 & 4 & first chunk of the control byte string (chunks are zero-indexed so it's always $0$ and should be removed in the future)\\
36--39 & 4 & how many control chunks (counts in how many chunks hash/start/count binary array will use)\\
40--43 & 4 & first chunk of the position array (chunks are zero-indexed so it's always $0$ and should be removed in the future)\\
44--47 & 4 & how many position chunks (counts in how many chunks position binary array will use)\\
48--55 & 8 & how many controls (number of items in each of hash/start/count arrays)\\
56--63 & 8 & how many positions (number of items in the position array)\\
64--71 & 8 & byte size of the control byte string\\
72--79 & 8 & byte size of the position array\\
80--87 & 8 & checksum of the control byte string\\
88--95 & 8 & checksum of the position array\\
96 & 1 & hash width: 2, 4, or 8\\
97 & 1 & start-array width: 2, 4, or 8\\
98 & 1 & count width: 2, 4, or 8\\
99 & 1 & position width: 2, 4, or 8\\
100 & 1 & Used to represent differnt orders of hash/start/count arrays. Now meaningless and fixed to $1$ \\
101 & 1 & Used to reresent how many positions there are per one hash entry. It's always $1$, so this is meaningless. \\
102 & 1 & Meaningless now, fixed to $2$ \\
103 & 1 & Number of extra entries in start aray at the end. Now it's meaningless and fixed to $1$ (and we don't need to store this one extra item at all)\\
104--127 & 24 & zeros\\
\bottomrule
\end{tabularx}
\label{tab:fields-in-grouped-mpf-directory-entry}
\end{table}

\subsubsection{Binary coding for MPS}

The MPS index stores the same rows as non-grouped MPF, but this time we also have additional one directory row.
Also the fragment keys are just not reffering to fragments anymore, but they represent groups of hashes similarily to minimap2~\ref{fig:index-mmi-layout}.

The directory contains metadata information like how many distinct hashes are there in total and for every group it stores min(hash), max(hash)
to allow us to quickly skip groups if certain hash values do not exist there.


\begin{figure}[htbp]
\centering
\begin{tikzpicture}[
  x=1cm,y=1cm,font=\footnotesize,
  lab/.style={align=center,font=\scriptsize,inner sep=1pt},
  cell/.style={align=center,font=\scriptsize,inner sep=1pt},
  zoom/.style={line width=0.85pt,line join=miter,line cap=rect,
    -{Latex[length=2.4mm,width=1.7mm]}}
]
  % ---------- (a) logical MPF ----------
  \node[anchor=south west,font=\bfseries\footnotesize] at (0.15,11.55)
    {(a) MPF: keep each source interval together};
  \draw[black!45,fill=black!3] (0.15,9.35) rectangle (4.55,11.30);
  \node[lab] at (2.35,10.95) {fragment 0: $[0,100)$};
  \fill[indexBlue!18] (0.45,9.60) rectangle (4.25,10.35);
  \draw[indexBlue] (0.45,9.60) rectangle (4.25,10.35);
  \node[cell] at (2.35,9.97) {$a$: 17, 65};

  \draw[black!45,fill=black!3] (4.85,9.35) rectangle (9.25,11.30);
  \node[lab] at (7.05,10.95) {fragment 1: $[100,200)$};
  \fill[indexOrange!20] (5.15,9.95) rectangle (8.95,10.55);
  \draw[indexOrange] (5.15,9.95) rectangle (8.95,10.55);
  \node[cell] at (7.05,10.25) {$b$: 120};
  \fill[indexBlue!18] (5.15,9.55) rectangle (8.95,9.90);
  \draw[indexBlue] (5.15,9.55) rectangle (8.95,9.90);
  \node[cell] at (7.05,9.72) {$a$: 180};

  \draw[black!45,fill=black!3] (9.55,9.35) rectangle (14.35,11.30);
  \node[lab] at (11.95,10.95) {fragment 2: $[200,300)$};
  \fill[indexBlue!18] (9.85,9.95) rectangle (14.05,10.55);
  \draw[indexBlue] (9.85,9.95) rectangle (14.05,10.55);
  \node[cell] at (11.95,10.25) {$a$: 230};
  \fill[indexGreen!20] (9.85,9.55) rectangle (14.05,9.90);
  \draw[indexGreen] (9.85,9.55) rectangle (14.05,9.90);
  \node[cell] at (11.95,9.72) {$c$: 280};

  % ---------- (b) logical MPS ----------
  \node[anchor=south west,font=\bfseries\footnotesize] at (0.15,8.85)
    {(b) MPS: put equal hashes together across the sequence};
  \fill[indexBlue!18] (0.15,7.95) rectangle (1.35,8.50);
  \draw[indexBlue] (0.15,7.95) rectangle (1.35,8.50);
  \node[cell] at (0.75,8.22) {$a$};
  \foreach \x/\p in {2.15/17,3.45/65,4.75/180,6.05/230} {
    \fill[indexBlue!18] (\x,7.95) rectangle (\x+1.15,8.50);
    \draw[indexBlue] (\x,7.95) rectangle (\x+1.15,8.50);
    \node[cell] at (\x+0.575,8.22) {$\p$};
  }
  \draw[indexBlue,line width=0.8pt,-{Latex[length=2.2mm]}] (1.35,8.22) -- (2.15,8.22);

  \fill[indexOrange!20] (0.15,7.20) rectangle (1.35,7.75);
  \draw[indexOrange] (0.15,7.20) rectangle (1.35,7.75);
  \node[cell] at (0.75,7.47) {$b$};
  \fill[indexOrange!20] (2.15,7.20) rectangle (3.30,7.75);
  \draw[indexOrange] (2.15,7.20) rectangle (3.30,7.75);
  \node[cell] at (2.725,7.47) {$120$};
  \draw[indexOrange,line width=0.8pt,-{Latex[length=2.2mm]}] (1.35,7.47) -- (2.15,7.47);

  \fill[indexGreen!20] (0.15,6.45) rectangle (1.35,7.00);
  \draw[indexGreen] (0.15,6.45) rectangle (1.35,7.00);
  \node[cell] at (0.75,6.72) {$c$};
  \fill[indexGreen!20] (2.15,6.45) rectangle (3.30,7.00);
  \draw[indexGreen] (2.15,6.45) rectangle (3.30,7.00);
  \node[cell] at (2.725,6.72) {$280$};
  \draw[indexGreen,line width=0.8pt,-{Latex[length=2.2mm]}] (1.35,6.72) -- (2.15,6.72);
  \node[anchor=west,align=left,font=\scriptsize] at (8.0,7.35)
    {Equal hashes become one list.\\A query for $a$ reads that list only.};

  % ---------- (c) ordinary MPF byte string ----------
  \node[anchor=south west,font=\bfseries\footnotesize] at (0.15,5.72)
    {(c) Ordinary MPF: one fragment is one byte string};
  \node[lab,anchor=east] at (3.25,5.15)
    {$\langle$query\_key,\\fragment$\rangle$};
  \fill[black!8] (3.35,4.80) rectangle (5.55,5.50);
  \draw[black!55] (3.35,4.80) rectangle (5.55,5.50);
  \node[lab] at (4.45,5.15) {header\\32\,B};
  \fill[indexBlue!12] (5.55,4.80) rectangle (9.85,5.50);
  \draw[black!55] (5.55,4.80) rectangle (9.85,5.50);
  \node[lab] at (7.70,5.15) {controls\\$n\times 16$\,B};
  \fill[indexOrange!16] (9.85,4.80) rectangle (14.25,5.50);
  \draw[black!55] (9.85,4.80) rectangle (14.25,5.50);
  \node[lab] at (12.05,5.15) {positions\\$m\times 8$\,B};

  % 32-byte header, left edge aligned with the header. Not a separate panel.
  \fill[black!8] (3.35,2.55) rectangle (4.63,3.20);
  \draw[black!55] (3.35,2.55) rectangle (4.63,3.20);
  \node[lab] at (3.99,2.875) {magic\\4\,B};
  \fill[black!8] (4.63,2.55) rectangle (5.91,3.20);
  \draw[black!55] (4.63,2.55) rectangle (5.91,3.20);
  \node[lab] at (5.27,2.875) {revision\\4\,B};
  \fill[black!6] (5.91,2.55) rectangle (8.47,3.20);
  \draw[black!55] (5.91,2.55) rectangle (8.47,3.20);
  \node[lab] at (7.19,2.875) {length\\8\,B};
  \fill[indexBlue!14] (8.47,2.55) rectangle (9.75,3.20);
  \draw[black!55] (8.47,2.55) rectangle (9.75,3.20);
  \node[lab] at (9.11,2.875) {$n$\\4\,B};
  \fill[black!6] (9.75,2.55) rectangle (11.03,3.20);
  \draw[black!55] (9.75,2.55) rectangle (11.03,3.20);
  \node[lab] at (10.39,2.875) {reserved\\4\,B};
  \fill[indexOrange!16] (11.03,2.55) rectangle (13.59,3.20);
  \draw[black!55] (11.03,2.55) rectangle (13.59,3.20);
  \node[lab] at (12.31,2.875) {$m$\\8\,B};
  \draw[zoom,black!70] (4.45,4.80) -- (4.45,3.20);

  % one control, directly under the controls rectangle of (c). Not a separate panel.
  \fill[indexBlue!18] (5.55,3.65) rectangle (7.70,4.35);
  \draw[indexBlue] (5.55,3.65) rectangle (7.70,4.35);
  \node[lab] at (6.625,4.00) {hash\\8\,B};
  \fill[indexOrange!20] (7.70,3.65) rectangle (8.775,4.35);
  \draw[indexOrange] (7.70,3.65) rectangle (8.775,4.35);
  \node[lab] at (8.24,4.00) {start\\4\,B};
  \fill[indexGreen!20] (8.775,3.65) rectangle (9.85,4.35);
  \draw[indexGreen] (8.775,3.65) rectangle (9.85,4.35);
  \node[lab] at (9.31,4.00) {count\\4\,B};
  \draw[zoom,indexBlue] (7.70,4.80) -- (7.70,4.35);

  % (d) grouped MPF. Separate byte strings. No arrow back to the control in (c).
  \node[anchor=south west,font=\bfseries\footnotesize] at (0.15,1.98)
    {(d) Grouped MPF};
  \node[lab,anchor=east] at (3.25,1.50)
    {$\langle$query\_key,\\sub\_key$\rangle$};
  \fill[indexBlue!16] (4.50,1.15) rectangle (7.75,1.85);
  \draw[indexBlue] (4.50,1.15) rectangle (7.75,1.85);
  \node[lab] at (6.125,1.50) {hash array\\$n\times$ 2, 4, or 8\,B};
  \fill[indexOrange!18] (7.75,1.15) rectangle (11.00,1.85);
  \draw[indexOrange] (7.75,1.15) rectangle (11.00,1.85);
  \node[lab] at (9.375,1.50) {start array\\$(n{+}1)\times$ 2, 4, or 8\,B};
  \fill[indexGreen!18] (11.00,1.15) rectangle (14.25,1.85);
  \draw[indexGreen] (11.00,1.15) rectangle (14.25,1.85);
  \node[lab] at (12.625,1.50) {count array\\$n\times$ 2, 4, or 8\,B};

  \node[lab,anchor=east] at (3.25,0.52)
    {$\langle$query\_key,\\$2^{61}{+}$sub\_key$\rangle$};
  \fill[indexOrange!14] (4.50,0.17) rectangle (14.25,0.87);
  \draw[indexOrange] (4.50,0.17) rectangle (14.25,0.87);
  \node[lab] at (9.375,0.52) {position array\\$m\times$ 2, 4, or 8\,B};

  \node[lab,anchor=east] at (3.25,-1.70)
    {$\langle$query\_key,\\$2^{63}{-}1\rangle$};
  \node[lab] at (3.70,-1.70) {$p\times$\\128\,B};
  \fill[black!6] (4.22,-2.20) rectangle (5.44,-1.20);
  \draw[black!55] (4.22,-2.20) rectangle (5.44,-1.20);
  \node[lab] at (4.83,-1.70) {partition\\4\,B};
  \fill[black!8] (5.44,-2.20) rectangle (6.66,-1.20);
  \draw[black!55] (5.44,-2.20) rectangle (6.66,-1.20);
  \node[lab] at (6.05,-1.70) {fragment\\4\,B};
  \fill[indexBlue!14] (6.66,-2.20) rectangle (8.08,-1.20);
  \draw[black!55] (6.66,-2.20) rectangle (8.08,-1.20);
  \node[lab] at (7.37,-1.70) {min(hash)\\8\,B};
  \fill[indexBlue!14] (8.08,-2.20) rectangle (9.56,-1.20);
  \draw[black!55] (8.08,-2.20) rectangle (9.56,-1.20);
  \node[lab] at (8.82,-1.70) {max(hash)\\8\,B};
  \fill[black!8] (9.56,-2.20) rectangle (10.71,-1.20);
  \draw[black!55] (9.56,-2.20) rectangle (10.71,-1.20);
  \node[lab] at (10.135,-1.70) {sub-key\\8\,B};
  \fill[indexGreen!12] (10.71,-2.20) rectangle (11.60,-1.20);
  \draw[black!55] (10.71,-2.20) rectangle (11.60,-1.20);
  \node[lab] at (11.155,-1.70) {4\,B};
  \fill[indexGreen!12] (11.60,-2.20) rectangle (12.49,-1.20);
  \draw[black!55] (11.60,-2.20) rectangle (12.49,-1.20);
  \node[lab] at (12.045,-1.70) {4\,B};
  \fill[indexOrange!14] (12.49,-2.20) rectangle (13.38,-1.20);
  \draw[black!55] (12.49,-2.20) rectangle (13.38,-1.20);
  \node[lab] at (12.935,-1.70) {4\,B};
  \fill[indexOrange!14] (13.38,-2.20) rectangle (14.25,-1.20);
  \draw[black!55] (13.38,-2.20) rectangle (14.25,-1.20);
  \node[lab] at (13.815,-1.70) {4\,B};
  \node[lab,anchor=south] at (11.155,-0.22) {controls first};
  \draw[black!60,line width=0.7pt] (11.155,-0.22) -- (11.155,-1.20);
  \node[lab,anchor=south] at (12.935,-0.62) {positions first};
  \draw[black!60,line width=0.7pt] (12.935,-0.62) -- (12.935,-1.20);
  \node[lab,anchor=north] at (12.045,-2.45) {count};
  \draw[black!60,line width=0.7pt] (12.045,-2.20) -- (12.045,-2.45);
  \node[lab,anchor=north] at (13.815,-2.45) {count};
  \draw[black!60,line width=0.7pt] (13.815,-2.20) -- (13.815,-2.45);

  % (e) MPS: same byte string, then its directory, header, and one entry
  \node[anchor=south west,font=\bfseries\footnotesize] at (0.15,-3.25)
    {(e) MPS group: the same byte string as (c)};
  \node[lab,anchor=east] at (3.25,-3.77)
    {$\langle$query\_key,\\group number$\rangle$};
  \fill[black!8] (3.35,-4.12) rectangle (5.55,-3.42);
  \draw[black!55] (3.35,-4.12) rectangle (5.55,-3.42);
  \node[lab] at (4.45,-3.77) {header\\32\,B};
  \fill[indexBlue!12] (5.55,-4.12) rectangle (9.85,-3.42);
  \draw[black!55] (5.55,-4.12) rectangle (9.85,-3.42);
  \node[lab] at (7.70,-3.77) {controls\\$n\times 16$\,B};
  \fill[indexOrange!16] (9.85,-4.12) rectangle (14.25,-3.42);
  \draw[black!55] (9.85,-4.12) rectangle (14.25,-3.42);
  \node[lab] at (12.05,-3.77) {positions\\$m\times 8$\,B};

  \node[lab,anchor=east] at (3.25,-4.815)
    {$\langle$query\_key,\\$2^{63}{-}1\rangle$};
  \fill[black!10] (3.35,-5.19) rectangle (6.15,-4.44);
  \draw[black!55] (3.35,-5.19) rectangle (6.15,-4.44);
  \node[lab] at (4.75,-4.815) {directory header\\40\,B};
  \fill[indexBlue!12] (6.15,-5.19) rectangle (14.25,-4.44);
  \draw[black!55] (6.15,-5.19) rectangle (14.25,-4.44);
  \node[lab] at (10.20,-4.815) {one entry per group, 24\,B};

  \fill[black!8] (3.35,-7.10) rectangle (3.94,-6.20);
  \draw[black!55] (3.35,-7.10) rectangle (3.94,-6.20);
  \node[lab] at (3.645,-6.65) {4\,B};
  \fill[black!8] (3.94,-7.10) rectangle (4.53,-6.20);
  \draw[black!55] (3.94,-7.10) rectangle (4.53,-6.20);
  \node[lab] at (4.235,-6.65) {4\,B};
  \fill[black!6] (4.53,-7.10) rectangle (5.12,-6.20);
  \draw[black!55] (4.53,-7.10) rectangle (5.12,-6.20);
  \node[lab] at (4.825,-6.65) {4\,B};
  \fill[black!6] (5.12,-7.10) rectangle (5.71,-6.20);
  \draw[black!55] (5.12,-7.10) rectangle (5.71,-6.20);
  \node[lab] at (5.415,-6.65) {4\,B};
  \fill[indexBlue!14] (5.71,-7.10) rectangle (6.88,-6.20);
  \draw[black!55] (5.71,-7.10) rectangle (6.88,-6.20);
  \node[lab] at (6.295,-6.65) {groups\\8\,B};
  \fill[indexBlue!14] (6.88,-7.10) rectangle (8.05,-6.20);
  \draw[black!55] (6.88,-7.10) rectangle (8.05,-6.20);
  \node[lab] at (7.465,-6.65) {distinct\\hashes\\8\,B};
  \fill[indexOrange!16] (8.05,-7.10) rectangle (9.22,-6.20);
  \draw[black!55] (8.05,-7.10) rectangle (9.22,-6.20);
  \node[lab] at (8.635,-6.65) {positions\\8\,B};
  \node[lab,anchor=south] at (3.645,-5.72) {magic};
  \draw[black!60,line width=0.7pt] (3.645,-5.72) -- (3.645,-6.20);
  \node[lab,anchor=south] at (4.825,-5.72) {flags};
  \draw[black!60,line width=0.7pt] (4.825,-5.72) -- (4.825,-6.20);
  \node[lab,anchor=north] at (4.235,-7.35) {revision};
  \draw[black!60,line width=0.7pt] (4.235,-7.10) -- (4.235,-7.35);
  \node[lab,anchor=north] at (5.415,-7.35) {reserved};
  \draw[black!60,line width=0.7pt] (5.415,-7.10) -- (5.415,-7.35);
  \draw[zoom,black!70] (4.25,-5.19) -- (4.25,-6.20);

  \fill[indexBlue!16] (9.40,-7.10) rectangle (10.95,-6.20);
  \draw[indexBlue] (9.40,-7.10) rectangle (10.95,-6.20);
  \node[lab] at (10.175,-6.65) {min(hash)\\8\,B};
  \fill[indexBlue!16] (10.95,-7.10) rectangle (12.50,-6.20);
  \draw[indexBlue] (10.95,-7.10) rectangle (12.50,-6.20);
  \node[lab] at (11.725,-6.65) {max(hash)\\8\,B};
  \fill[black!8] (12.50,-7.10) rectangle (14.05,-6.20);
  \draw[black!55] (12.50,-7.10) rectangle (14.05,-6.20);
  \node[lab] at (13.275,-6.65) {group number\\8\,B};
  \draw[zoom,indexBlue] (10.20,-5.19) -- (10.20,-6.20);
\end{tikzpicture}
\caption[Byte strings of ordinary MPF, grouped MPF, and MPS]{
  Each row with two key in brackets $\langle$query\_key, sub\_key$\rangle$ represent binary string that will be split into multiple chunks that could be find 
  by complex key $\langle$query\_key, sub\_key, chunk\_number$\rangle$.
  There's always one special sub-key for directory for grouped MPF and MPS, these formats also store two rows per one sub-key.
}
\label{fig:index-mpf-mps}
\end{figure}

MPS index building needs to spend additional work on sorting and merging so later queries can skip
unrelated hash ranges, while MPF just keeps smaller hashmaps based on fragments. 
The current default BioseqDB active minimizer index is MPF.

\subsection{Index versioning}
\label{sec:index-compatibility}

Internal sequence storage tables use \path{epoch} to provide sequence versioning~\ref{lst:storage-sequence-locator}, indexes don't have to use separate versioning field.
I decided not to store epoch in the index table~\ref{tab:fields-in-grouped-mpf-directory-entry}, but instead \path{query_key} is number of the row inside the 
table shard. They are stable with respect to \path{INSERT}s, \path{DELETE}s and \path{UPDATE}s as update will create new sequence version. 
Special handling is only required when tables are rewritten using \path{VACCUM FULL}: vacuum full would cause the rows to change numbering, however
the indexes tables contain a hash of database, table space and Postgres \path{relfilenode}, so if the table is rewritten (\path{relfilenode} will change) BisoeqDB won't detect the old index.
This is the most naive approach, but provides basic protection against table rewrites.


\subsection{Symbol count index}
\label{sec:index-symbol-counts}

I implemented one index type that does not store minimizers at all: \texttt{symbol\_counts}
index keeps a very small symbol histogram for one sequence.
Each symbol has 64-bit counter (for DNA4, the counters are $A,C,G,T$ so sequence
\texttt{AACGTA} gives $(3,1,1,1)$). The counts are stored separately because as we use heap rows the value will be stored inline in Postgres table.
This allows very fast access and usually users don't need to count symbols.

\begin{figure}[htbp]
\centering
\begin{tikzpicture}[x=1cm,y=.7cm,font=\footnotesize]
  \node[anchor=west,font=\bfseries] at (0,5.6) {One sequence};
  \foreach \x/\s in {1/A,2/A,3/C,4/G,5/T,6/A}
    \node[index blue] at (\x,4.7) {\texttt{\s}};
  \node[anchor=west,font=\bfseries] at (8,5.6) {Four counters};
  \foreach \x/\s/\n in {8.5/A/3,9.7/C/1,10.9/G/1,12.1/T/1} {
    \draw[fill=indexGreen!18,draw=indexGreen] (\x-.38,.7) rectangle (\x+.38,.7+\n);
    \node[above] at (\x,.7+\n) {$\n$};
    \node[below] at (\x,.7) {\texttt{\s}};
  }
  \draw[index arrow] (6.5,3.6)--(7.75,3.6);
  \node[anchor=west,align=left] at (0,1.6)
    {To count A, read A's counter.\\Seed lookup: positions are still needed.};
\end{tikzpicture}
\caption{Hitsogram counter for symbols is stored separately so BioseqDB access it only when specified in the query}
\label{fig:index-symbol-counts}
\end{figure}

\begin{lstlisting}[style=bioseq-index-sql,
  caption={Request a symbol count for each stored sequence.},
  label={lst:index-symbol-counts}]
SELECT id, bioseq.symbol_count(seq, 'A') AS a_count
FROM reads;
\end{lstlisting}
Using this small metadata index allows me to never access the sequence content at all, meaning the read is very fast.
If the index entry is not found, then BioseqDB will traverse entire sequence.
The primary benefit is that this can be computed when loading FASTA or inserting rows autmatically so the user pays the 
cost only once (or later when infrequent \path{INSERT} or \path{UPDATE} happen)

\section{Building MPS}
\label{sec:index-construction}

Building MPF index is rather simple as we traverse each fragment and can separately construct hashmaps with minimizers (and hashmaps for fragmenst are relatively
small compared to all RAM capacity).
Constructing MPS is more complex as I need to somehow store all minimizer globally and this can potentially be bigger than entire RAM capacity.
Index builder could collect every $(hash,position,strand)$ record and
sort the complete vector, but for maps bigger than RAM this couldn't be done efficiently, instead current code does similar thing to MPF with a slight differences:
\begin{itemize}
  \item To reduce memory access costs as much as possible MPF treats each fragment as entirely separate sequence, thus minimizers on edges won't be generated (the first and last $k+w-2$ symbols). 
  This reduces minimizer density at edges but means that we access only one fragment to generate index entry.
  \item For MPS I use subtly different approach and I collect all minimizers. This means that during parallelized fragment read, the reader need to read not only own fragment, but also 
  parts of fragment to the right and left. This is depicted as "context" on the figure below~\ref{fig:index-build-halos}.
\end{itemize}

\begin{figure}[htbp]
\centering
\begin{tikzpicture}[x=1cm,y=1cm,font=\footnotesize]
  \node[anchor=west,font=\bfseries] at (0,5.1) {Neighboring construction cores};
  \draw[fill=indexBlue!12,draw=indexBlue] (1,3.25) rectangle (6,3.85);
  \draw[fill=indexGreen!12,draw=indexGreen] (6,3.25) rectangle (11,3.85);
  \node at (3.5,3.55) {fragment 0}; \node at (8.5,3.55) {fragment 1};
  \draw[dashed,black!55] (6,.25)--(6,4.15);
  \node[above] at (6,4.15) {fragment 0 ends at positon $p$ of sequence};
  \draw[fill=indexBlue!8,draw=indexBlue] (1,2.15) rectangle (7.2,2.62);
  \draw[pattern=north east lines,pattern color=indexBlue,draw=indexBlue]
    (6,2.15) rectangle (7.2,2.62);
  \draw[fill=indexGreen!8,draw=indexGreen] (4.8,1.05) rectangle (11,1.52);
  \draw[pattern=north east lines,pattern color=indexGreen,draw=indexGreen]
    (4.8,1.05) rectangle (6,1.52);
  \node[anchor=east] at (.8,2.38) {read 0};
  \node[anchor=east] at (.8,1.28) {read 1};
  \node[anchor=west] at (7.45,2.38) {right context};
  \node[anchor=east] at (4.55,1.28) {left context};
  \draw[very thick,indexOrange] (5.25,.5)--(6.25,.5);
  \fill[indexOrange] (6.25,.5) circle[radius=2.5pt];
  \node[anchor=west] at (6.5,.5) {seed ending in fragment 1: emitted by fragment 1};
\end{tikzpicture}
\caption{Each reader to generate minimziers for one fragment in MPS read also neighbouring fragments so first and last minimizers in fragment
can be generated. Size of "context" is $w+k-2$ symbols.}
\label{fig:index-build-halos}
\end{figure}

\begin{minipage}{\linewidth}
MPS does the following steps to build entire index:
\begin{enumerate}
  \item Read fragments and for one fragment compute minimizers, sort them by hash and remove duplicates similarily to MPF.
  \item Write this computed MPF-like part into temporary file storage.
  \item After all fragments have temporary files, then select fixed number of them and merge them into one sorted temporary file.
  \item Continue this merging procedure until one temporary file remains.
  \item Read the final temporary file and build directory row inside storage.
\end{enumerate}
\end{minipage}

\begin{figure}[htbp]
\centering
\begin{tikzpicture}[x=1cm,y=1cm,font=\footnotesize]
  \node[anchor=west,font=\bfseries] at (0,5.8) {Small sorted runs};
  \foreach \x/\name/\items in {1.4/0/{a,c,f},4.5/1/{a,b,e},7.6/2/{b,d,f},10.7/3/{a,c,e}} {
    \node[index blue,minimum width=2.5cm] (r\name) at (\x,4.8) {run \name\\$\items$};
  }
  \node[index orange,minimum width=3.6cm] (m0) at (3,3.0) {$a,a,b,c,e,f$};
  \node[index orange,minimum width=3.6cm] (m1) at (9,3.0) {$a,b,c,d,e,f$};
  \draw[index arrow] (r0.south)--(m0.north west);
  \draw[index arrow] (r1.south)--(m0.north east);
  \draw[index arrow] (r2.south)--(m1.north west);
  \draw[index arrow] (r3.south)--(m1.north east);
  \node[anchor=west,align=left] at (0,1.7)
    {Merging fixed count per iteration - in this case $2$ runs at the time};
  \node[index green,minimum width=6.8cm] (final) at (7.6,1.2)
    {$a,a,a,b,b,c,c,d,e,e,f,f$};
  \draw[index arrow] (m0.south)--(final.north west);
  \draw[index arrow] (m1.south)--(final.north east);
  \node[anchor=west,text=black!65] at (0,.1)
    {Each entry can be very large and we don't have guarantee that any merged run will fit in RAM};
\end{tikzpicture}
\caption{Merging of hash entries (hash value and minimizer position)until one final entry is left and can be stored permamently in index storage. Letters represent different hash values.}
\label{fig:index-mps-merge}
\end{figure}

The spill path uses PostgreSQL's \texttt{BufFile} temporary-file interface.
The sorter retains sorted runs in RAM when memory admission succeeds and
spills them when it cannot reserve that memory. It can also move a retained
run to a temporary file to make room for later work. Thus the merge structure
in Figure~\ref{fig:index-mps-merge} does not require every run to be written
to disk.

\section{Selective reads}
\label{sec:index-cold-loading}
\label{sec:index-mps-selective-reads}

Main benefit of MPS is that each group can be a single row within Postgres
(if $n$, $m$ are not causing the array to spill over multiple chunks~\ref{fig:index-storage-layouts}).
This means that minimap2 has to read and check every entry in index, while BioseqDB can only read entiore directory row 
and using $min(hash)$ and $max(hash)$ values to decide if we should read certain groups of minimizers.

The similar logic applies to grouped MPF index as well, but in case of grouped MPF we can also check count array and if the count is very high, then 
we can abandon readind reparate position array at all (the most frequent minimizers should be filtered out as they are not good candiates for anchors). 

The hash-major MPS reader performs the following steps:
\begin{enumerate}
  \item Read and check the complete directory, then find hash ranges that intersect with query hashes.
  \item Read the selected controls to find each key's exact global occurrence count
  \item Apply the occurence count threshold: if the count is high, we don't process this hash
  \item Read the chunks selected by this filtering and read positions form them
\end{enumerate}

\begin{figure}[htbp]
\centering
\begin{tikzpicture}[x=1cm,y=1cm,font=\footnotesize]
  \node[anchor=west,font=\bfseries] at (0,6.25) {Query hashes $\{12,35,78\}$; occurrence limit 5};
  \node[anchor=east] at (1.2,5.35) {directory};
  \foreach \x/\range/\col in {2.6/{0--19}/indexBlue,5.2/{20--39}/indexBlue,7.8/{40--59}/black,10.4/{60--79}/indexBlue}
    \node[index cell,fill=\col!10,draw=\col,minimum width=2.3cm] at (\x,5.35) {$\range$};
  \node[anchor=east] at (1.2,3.95) {controls};
  \node[index green,minimum width=2.3cm] (c12) at (2.6,3.95) {$12:\ count=2$};
  \node[index orange,minimum width=2.3cm] (c35) at (5.2,3.95) {$35:\ count=90$};
  \node[index cell,dashed,minimum width=2.3cm] at (7.8,3.95) {not read};
  \node[index green,minimum width=2.3cm] (c78) at (10.4,3.95) {$78:\ count=3$};
  \foreach \x in {2.6,5.2,10.4}
    \draw[index arrow] (\x,5.04)--(\x,4.28);
  \node[anchor=east] at (1.2,2.0) {postings};
  \draw[fill=indexGreen!16,draw=indexGreen] (1.45,1.65) rectangle (3.1,2.35);
  \node at (2.275,2.0) {$17,42$};
  \draw[fill=black!7,draw=black!45] (3.25,1.65) rectangle (8.9,2.35);
  \node at (6.075,2.0) {90 positions: not read};
  \draw[fill=indexGreen!16,draw=indexGreen] (9.05,1.65) rectangle (11.55,2.35);
  \node at (10.3,2.0) {$18,54,96$};
  \draw[index arrow,indexGreen] (c12.south)--(2.275,2.39);
  \draw[index arrow,indexGreen] (c78.south)--(10.3,2.39);
  \draw[thick,indexOrange] (5.2,3.61)--(5.2,2.92);
  \draw[thick,indexOrange] (5.02,2.92)--(5.38,2.92);
  \node[anchor=west,text=indexOrange] at (5.52,3.03) {$90>5$: stop before positions};
  \node[anchor=west,align=left] at (1.45,.6)
    {Query filters out the most frequent minimizers according to configured threshold and skips min-max ranges of hashes};
\end{tikzpicture}
\caption{Two decisions save different reads. All directory entries are
read; blue ranges select the controls. Counts reject a frequent hash before
its long list is loaded. Positions and thresholds are illustrative.}
\label{fig:index-selective-load}
\label{fig:format6-streams}
\end{figure}


\section{Other index considerations}
\label{sec:index-operational-considerations}

Indexes can automatically update themselves on index, however currently enabled by default is grouped MPF index, which due to bugs cannot correctly 
return minimizers which leads to recomputing them when no explicit index read policy is specified and it doesn't support update on insert.
More details are presented in the Appendix 1, where I also present concrete examples of how the BioseqDB should be used.

This bug was caused by the fact that grouped MPF not only produces relational index, but also saves that compressed relation chunks into separate file that is used to enable faster reads during JOINs.
The reader by default rejects ordinary single table queries, so you need to change the GUC to disable grouped-MPF (and use regular MPF instead).
Saving file enables the JOIN to concurrently access stored data which is a cheap workaround around the decision to only use single Postgres reader process.

\begin{figure}[H]
\centering
\begin{tikzpicture}[x=1cm,y=1cm,font=\footnotesize]
  \foreach \shift/\mode/\heading in
    {0/0/{(a) Ordinary MPF},7/1/{(b) Grouped MPF}} {
    \begin{scope}[xshift=\shift cm]
      \node[anchor=west,font=\bfseries] at (.1,10.65) {\heading};
      \node at (4.1,10.15) {\texttt{refs}: targets};
      \draw[fill=black!4,draw=black!55] (2.4,8.08) rectangle (5.8,9.88);
      \draw[fill=indexGreen!16,draw=indexGreen] (2.4,8.68) rectangle (5.8,9.28);
      \draw[black!55] (2.4,9.28)--(5.8,9.28);
      \draw[black!55] (3.45,8.08)--(3.45,9.88);
      \node at (2.925,9.58) {id};
      \node at (4.625,9.58) {seq};
      \node at (2.925,8.98) {$t_0$};
      \node at (4.625,8.98) {$\cdots$};
      \node[text=black!50] at (2.925,8.38) {$t_1$};
      \node[text=black!50] at (4.625,8.38) {$\cdots$};
      \draw[index arrow] (4.1,7.98)--(4.1,7.27);
      \node[anchor=west] at (4.3,7.68) {index for $t_0$};

      \ifnum\mode=0
        \node at (4.1,7.04) {MPF index relation};
        \foreach \yy in {6.6,6.23,5.86} {
          \draw[fill=black!6,draw=black!55] (2.7,\yy-.12) rectangle (5.5,\yy+.12);
          \draw[fill=indexOrange!25,draw=indexOrange] (3.25,\yy-.12) rectangle (5.5,\yy+.12);
        }
        \draw[index arrow,indexOrange] (4.1,5.59)--(4.1,4.75);
        \node[anchor=west,align=left] at (4.32,5.18) {load and\\assemble};
        \draw[fill=indexOrange!10,draw=indexOrange] (3.0,2.25) rectangle (6.15,4.65);
        \node at (4.575,4.3) {private CPU lookup};
      \else
        \node at (4.1,7.04) {MPF reader image};
        \draw[fill=indexGreen!7,draw=indexGreen]
          (2.7,5.74)--(5.5,5.74)--(5.5,6.43)--(5.21,6.72)--(2.7,6.72)--cycle;
        \draw[indexGreen] (5.21,6.72)--(5.21,6.43)--(5.5,6.43);
        \foreach \yy in {6.47,6.2,5.93}
          \draw[line width=3pt,indexGreen!45] (2.98,\yy)--(4.92,\yy);
        \draw[index arrow,indexGreen] (4.1,5.59)--(4.1,4.75);
        \node[anchor=west,align=left] at (4.32,5.18) {open and\\\texttt{mmap}};
        \draw[dashed,fill=indexGreen!6,draw=indexGreen] (3.0,2.25) rectangle (6.15,4.65);
        \node at (4.575,4.3) {mapped array view};
      \fi
      \node[align=center] at (4.575,3.36) {one reader\\for target $t_0$};

      \node at (1.0,5.02) {\texttt{queries}};
      \draw[fill=indexBlue!6,draw=indexBlue] (.1,2.25) rectangle (1.9,4.65);
      \draw[indexBlue] (.75,2.25)--(.75,4.65);
      \foreach \yy in {4.05,3.45,2.85}
        \draw[indexBlue] (.1,\yy)--(1.9,\yy);
      \node at (.425,4.35) {id};
      \node at (1.325,4.35) {seq};
      \foreach \yy/\q in {3.75/0,3.15/1,2.55/2} {
        \node at (.425,\yy) {$q_{\q}$};
        \node at (1.325,\yy) {$\cdots$};
        \draw[index arrow,indexBlue] (1.98,\yy)--(2.9,\yy);
      }
      \node[align=center] at (1.0,1.72) {keep query\\rows};
      \draw[index arrow] (4.575,2.14)--(4.575,1.55);
      \node at (4.575,1.27) {target positions for anchors};
      \draw[black!45] (3.1,.8)--(6.05,.8);
      \foreach \xx in {3.32,4.02,5.7}
        \fill[indexGreen] (\xx,.8) circle (2.2pt);
    \end{scope}
  }
\end{tikzpicture}
\caption{For large joins grouped-MPF uses saved direct binary dump of relational index, which speeds up the alignment by reducing time to sequentially read the backing table}
\label{fig:index-join-read-paths}
\end{figure}

All indexes in BioseqDB can be automatically updated on insert (but as noted earlier, grouped-MPF will produce error on insert due to buggy implementation)
\begin{table}[htbp]
\centering
\small
\caption{Maintenance policies}
\label{tab:index-update-policy}
\begin{tabularx}{\linewidth}{@{}p{.31\linewidth}X@{}}
\toprule
\textbf{Policy} & \textbf{Description}\\
\midrule
\texttt{UPDATE\_ON\_INSERT} & When user inserts data, we automatically compute and store minimizers for new sequence \\
\texttt{UPDATE\_MANUALLY} & User needs to manually call \texttt{bioseq.update\_index} \\
\bottomrule
\end{tabularx}
\end{table}

\begin{table}[htbp]
\centering
\small
\caption{Default policies in BioseqDB}
\label{tab:index-policy-defaults}
\begin{tabularx}{\linewidth}{@{}p{.34\linewidth}p{.30\linewidth}X@{}}
\toprule
\textbf{Index} & \textbf{Policy for user updates} & \textbf{Description}\\
\midrule
\path{minimizers_per_fragment} & \texttt{UPDATE\_ON\_INSERT} & Default minimizer index\\
\path{minimizers_per_sequence} & \texttt{UPDATE\_MANUALLY} & Not selected by default\\
\path{symbol_counts} & \texttt{UPDATE\_ON\_INSERT} & Used when user demans to count symbols in a sequence\\
\bottomrule
\end{tabularx}
\end{table}


\flushbottom
